% GENERATED FILE. Edit canonical lessons, then run npm run generate:books.
% canonical-source-hash: fnv1a64:9ae6e258e0e1db05
% canonical-lessons: SA-C333-sapah, SA-C333-devah, SA-C333-varah, SA-C333-prasidati, SA-C333-amarah

\chapter{The Curse and the Boon}
\label{ch:sa-the-curse-and-the-boon}
\hlchaptermodality{333}

\begin{chapteropening}
\textbf{By the end of this chapter:} \emph{I can tell how a curse was laid and what followed.}

Complete the last lesson of chapter 333: I can tell how a curse was laid and what followed.
\end{chapteropening}

\section[\texorpdfstring{\sk{शापः} (śāpaḥ)}{शापः (śāpaḥ)}]{\sk{शापः} (śāpaḥ) — a curse}

\label{lesson:SA-C333-sapah}



\begin{quote}
Before the new one: say the Sanskrit for penance, then the Sanskrit for curses.
\end{quote}

\subsection*{You'll want to know: \sk{शापः}}
\textbf{\sk{शापः}} — \emph{śāpaḥ} — \textquotedblleft{}a curse\textquotedblright{}.

\begin{quote}
\sk{ऋषेः} \sk{शापेन} \sk{राजा} \sk{दुःखी} \sk{भवति।}
\end{quote}

\begin{quote}
\emph{ṛṣeḥ śāpena rājā duḥkhī bhavati.}
\end{quote}

\begin{quote}
The seer's curse makes the king unhappy.
\end{quote}

A masculine noun.

\begin{cousinweb}
\textbf{\sk{शापः}} (\emph{śāpaḥ}) comes from the root √śap, \textquotedblleft{}to curse\textquotedblright{}, as in \emph{śapati}; MW gives \emph{śāpa} as a curse, malediction, oath (MW).
\end{cousinweb}

\begin{grammarlens}[title={What you've built}]
One more word for telling what happened.
\end{grammarlens}

\subsection*{Your turn}
\begin{itemize}
\raggedright
  \item \emph{Say it:} \emph{śāpaḥ}
  \item \emph{Say it:} \emph{śāpaḥ}, once more
  \item \emph{Say it:} \emph{ṛṣeḥ śāpena rājā duḥkhī bhavati.} — The seer's curse makes the king unhappy.
  \item \emph{Recall:} say the Sanskrit for penance, then the Sanskrit for curses, then say \emph{śāpaḥ} again
\end{itemize}

\begin{tcolorbox}[breakable,colback=teal!4,colframe=teal!35!black,title={Before you move on}]
What does \sk{शापः} mean? (\textquotedblleft{}A curse\textquotedblright{}.) Say it once more.
\end{tcolorbox}
\section[\texorpdfstring{\sk{देवः} (devaḥ)}{देवः (devaḥ)}]{\sk{देवः} (devaḥ) — a god}

\label{lesson:SA-C333-devah}



\begin{quote}
Before the new one: say the Sanskrit for curses, then the Sanskrit for a curse.
\end{quote}

\subsection*{You'll want to know: \sk{देवः}}
\textbf{\sk{देवः}} — \emph{devaḥ} — \textquotedblleft{}a god\textquotedblright{}.

\begin{quote}
\sk{देवः} \sk{ऋषेः} \sk{तपः} \sk{पश्यति।}
\end{quote}

\begin{quote}
\emph{devaḥ ṛṣeḥ tapaḥ paśyati.}
\end{quote}

\begin{quote}
The god sees the seer's penance.
\end{quote}

A masculine noun.

\begin{cousinweb}
\textbf{\sk{देवः}} (\emph{devaḥ}) comes from the root √div, \textquotedblleft{}to shine\textquotedblright{}; MW gives \emph{deva} as heavenly, divine and, as a noun, a deity, a god (MW).
\end{cousinweb}

\begin{grammarlens}[title={What you've built}]
One more word for telling what happened.
\end{grammarlens}

\subsection*{Your turn}
\begin{itemize}
\raggedright
  \item \emph{Say it:} \emph{devaḥ}
  \item \emph{Say it:} \emph{devaḥ}, once more
  \item \emph{Say it:} \emph{devaḥ ṛṣeḥ tapaḥ paśyati.} — The god sees the seer's penance.
  \item \emph{Recall:} say the Sanskrit for curses, then the Sanskrit for a curse, then say \emph{devaḥ} again
\end{itemize}

\begin{tcolorbox}[breakable,colback=teal!4,colframe=teal!35!black,title={Before you move on}]
What does \sk{देवः} mean? (\textquotedblleft{}A god\textquotedblright{}.) Say it once more.
\end{tcolorbox}
\section[\texorpdfstring{\sk{वरः} (varaḥ)}{वरः (varaḥ)}]{\sk{वरः} (varaḥ) — a boon}

\label{lesson:SA-C333-varah}



\begin{quote}
Before the new one: say the Sanskrit for a curse, then the Sanskrit for a god.
\end{quote}

\subsection*{You'll want to know: \sk{वरः}}
\textbf{\sk{वरः}} — \emph{varaḥ} — \textquotedblleft{}a boon\textquotedblright{}.

\begin{quote}
\sk{देवः} \sk{ऋषये} \sk{वरं} \sk{ददाति।}
\end{quote}

\begin{quote}
\emph{devaḥ ṛṣaye varaṃ dadāti.}
\end{quote}

\begin{quote}
The god gives the seer a boon.
\end{quote}

A masculine noun.

\begin{cousinweb}
\textbf{\sk{वरः}} (\emph{varaḥ}) comes from the root √vṛ, \textquotedblleft{}to choose\textquotedblright{}, like the \emph{vara} who is chosen as a bridegroom; MW gives \emph{vara} as choosing, a wish, and a boon or gift (MW).
\end{cousinweb}

\begin{grammarlens}[title={What you've built}]
One more word for telling what happened.
\end{grammarlens}

\subsection*{Your turn}
\begin{itemize}
\raggedright
  \item \emph{Say it:} \emph{varaḥ}
  \item \emph{Say it:} \emph{varaḥ}, once more
  \item \emph{Say it:} \emph{devaḥ ṛṣaye varaṃ dadāti.} — The god gives the seer a boon.
  \item \emph{Recall:} say the Sanskrit for a curse, then the Sanskrit for a god, then say \emph{varaḥ} again
\end{itemize}

\begin{tcolorbox}[breakable,colback=teal!4,colframe=teal!35!black,title={Before you move on}]
What does \sk{वरः} mean? (\textquotedblleft{}A boon\textquotedblright{}.) Say it once more.
\end{tcolorbox}
\section[\texorpdfstring{\sk{प्रसीदति} (prasīdati)}{प्रसीदति (prasīdati)}]{\sk{प्रसीदति} (prasīdati) — shows favour; is gracious}

\label{lesson:SA-C333-prasidati}



\begin{quote}
Before the new one: say the Sanskrit for a god, then the Sanskrit for a boon.
\end{quote}

\subsection*{You'll want to know: \sk{प्रसीदति}}
\textbf{\sk{प्रसीदति}} — \emph{prasīdati} — \textquotedblleft{}shows favour; is gracious\textquotedblright{}.

\begin{quote}
\sk{तपः} \sk{दृष्ट्वा} \sk{देवः} \sk{प्रसीदति।}
\end{quote}

\begin{quote}
\emph{tapaḥ dṛṣṭvā devaḥ prasīdati.}
\end{quote}

\begin{quote}
Seeing the penance, the god shows favour.
\end{quote}

Past (he …): \sk{सः} \sk{अन्ते} \sk{प्रसन्नवान्}, \emph{saḥ ante prasannavān}.

\begin{cousinweb}
\textbf{\sk{प्रसीदति}} (\emph{prasīdati}) is built from \emph{pra} and the root √sad, \textquotedblleft{}to sit, settle\textquotedblright{}; MW gives \emph{pra-sad} as to settle down, grow clear and calm, and to be gracious (MW).
\end{cousinweb}

\begin{grammarlens}[title={What you've built}]
One more word for telling what happened.
\end{grammarlens}

\subsection*{Your turn}
\begin{itemize}
\raggedright
  \item \emph{Say it:} \emph{prasīdati}
  \item \emph{Say it:} \emph{prasīdati}, once more
  \item \emph{Say it:} \emph{tapaḥ dṛṣṭvā devaḥ prasīdati.} — Seeing the penance, the god shows favour.
  \item \emph{Recall:} say the Sanskrit for a god, then the Sanskrit for a boon, then say \emph{prasīdati} again
\end{itemize}

\begin{tcolorbox}[breakable,colback=teal!4,colframe=teal!35!black,title={Before you move on}]
What does \sk{प्रसीदति} mean? (\textquotedblleft{}Shows favour; is gracious\textquotedblright{}.) Say it once more.
\end{tcolorbox}
\section[\texorpdfstring{\sk{अमरः} (amaraḥ)}{अमरः (amaraḥ)}]{\sk{अमरः} (amaraḥ) — immortal}

\label{lesson:SA-C333-amarah}



\begin{quote}
Before the new one: say the Sanskrit for a boon, then the Sanskrit for shows favour.
\end{quote}

\subsection*{You'll want to know: \sk{अमरः}}
\textbf{\sk{अमरः}} — \emph{amaraḥ} — \textquotedblleft{}immortal\textquotedblright{}.

\begin{quote}
\sk{वरेण} \sk{ऋषिः} \sk{अमरः} \sk{भवति।}
\end{quote}

\begin{quote}
\emph{vareṇa ṛṣiḥ amaraḥ bhavati.}
\end{quote}

\begin{quote}
Through the boon the seer becomes immortal.
\end{quote}

\begin{cousinweb}
\textbf{\sk{अमरः}} (\emph{amaraḥ}) is built from \emph{a}, \textquotedblleft{}not\textquotedblright{}, and \emph{mara}, \textquotedblleft{}dying\textquotedblright{}, from the root √mṛ, as in \emph{mriyate}; MW gives \emph{amara} as undying, immortal, and as a noun a god (MW).
\end{cousinweb}

\begin{grammarlens}[title={What you've built}]
One more word for telling what happened.
\end{grammarlens}

\subsection*{Your turn}
\begin{itemize}
\raggedright
  \item \emph{Say it:} \emph{amaraḥ}
  \item \emph{Say it:} \emph{amaraḥ}, once more
  \item \emph{Say it:} \emph{vareṇa ṛṣiḥ amaraḥ bhavati.} — Through the boon the seer becomes immortal.
  \item \emph{Recall:} say the Sanskrit for a boon, then the Sanskrit for shows favour, then say \emph{amaraḥ} again
\end{itemize}

\begin{tcolorbox}[breakable,colback=teal!4,colframe=teal!35!black,title={Before you move on}]
What does \sk{अमरः} mean? (\textquotedblleft{}Immortal\textquotedblright{}.) Say it once more.
\end{tcolorbox}
